\subsection{拓扑向量空间的定义}

定义 1. --- 给定一个拓扑除环 K (GT, III, § 6.7) 和一个集合 E，使得 E 具有

\(1^{\circ}\) \(\mathbf{K}\) 上的左向量空间结构；

\(2^{\mathrm{o}}\) 一个与 E 的加法群结构相容的拓扑 (GT, III,

§ 1.1)，且此外还满足下面的公理：

(EVT) 映射 \((\lambda, x) \mapsto \lambda x\)（\(\mathbf{K} \times \mathbf{E}\) 到 \(\mathbf{E}\) 中的映射）连续，

则称 E 为 K 上的左拓扑向量空间。

这等价于说 E 是一个拓扑左 K-模 (GT, III, § 6.6)。

集合 E 上相对于 K 的左向量空间结构与 E 上给定的拓扑称为相容的，如果该拓扑与 E 的加法群结构相容，并且此外公理 (EVT) 成立。这也就是说，\((x, y) \mapsto x + y\) 与 \((\lambda, x) \mapsto \lambda x\) 这两个映射（分别为 \(E \times E\) 与 \(K \times E\) 到 E 中的映射）都连续；因为这时映射 \(x \mapsto -x = (-1)x\) 连续，从而 E 的拓扑与其加法群结构相容。

若 \(\mathbf{E}\) 是 \(\mathbf{K}\) 上的左拓扑向量空间，则我们称仅赋有向量空间结构的 \(\mathbf{E}\) 是拓扑向量空间 \(\mathbf{E}\) 的底空间。

与 R 上 E 的向量空间结构相容。因为，设 \(u_{0}\) 为 E 中一点，设 H 为 X 的紧子集，并设 \(\varepsilon\) 为任意一个严格正的数。实值函数 \(u_{0}\) 在 H 上有界；令 \(a = \sup_{t \in H} |u_{0}(t)|\)；若 u 是 E 的任意一点，则对一切 \(t \in H\)

\[
\left| \lambda u (t) - \lambda_ {0} u _ {0} (t) \right| \leqslant | \lambda |. \left| u (t) - u _ {0} (t) \right| + a \left| \lambda - \lambda_ {0} \right|.
\]

于是，若 \(|\lambda - \lambda_0| \leqslant \varepsilon\) 且 \(|u(t) - u_0(t)| \leqslant \varepsilon\) 对一切 \(t \in \mathbf{H}\) 成立，则对 \(t \in \mathbf{H}\) 有 \(|\lambda u(t) - \lambda_0 u_0(t)| \leqslant \varepsilon (\varepsilon + |\lambda_0| + a)\)，这表明公理 (EVT) 得到满足；同样可以验证，紧收敛拓扑与 \(\mathbf{E}\) 的加法群结构相容。

另一方面，若 X 不是紧的，（X 上的）一致收敛拓扑未必与 E 的向量空间结构相容；例如，若 X = R 且 \(u_{0}\) 是 R 上的一个无界连续函数，则 R 到 E 中的映射 \(\lambda \mapsto \lambda u_{0}\) 在 E 的一致收敛拓扑下不连续。

\begin{enumerate}
\def\labelenumi{\arabic{enumi})}
\setcounter{enumi}{4}
\tightlist
\item
  设 E 是拓扑除环 K 上有限维 n 的向量空间；存在一个向量 K-空间的同构 \(u:K_{s}^{n}\to E\)，而且，若 v 是 \(K_{s}^{n}\) 到 E 上的第二个同构，则可以写 \(v=u\circ f\)，其中 f 是向量 K-空间 \(K_{s}^{n}\) 的一个自同构。在 \(K_{s}^{n}\) 上考虑与其向量空间结构相容的乘积拓扑（例 3）；由于 \(K_{s}^{n}\) 到自身的每个线性映射对该拓扑都连续，故向量空间 \(K_{s}^{n}\) 的每个自同构都是双连续的。因此，若借助 \(K_{s}^{n}\) 到 E 上的随便哪一个同构，把 \(K_{s}^{n}\) 的乘积拓扑搬到 E 上，则如此在 E 上得到的拓扑与所用的是哪个同构无关；我们称之为 E 上的典范拓扑；当 K 是一个非离散的完备赋值除环时，我们将另行刻画它 (I, § 1.3)。对于 E 上的典范拓扑，E 到 K 上任一拓扑向量空间中的每个线性映射都连续。
\end{enumerate}

按照与定义 1 相同的方式，可以定义拓扑除环 K 上的右拓扑向量空间；但 K 上的每个右向量空间都可以看作 K 的反除环 \(K^{0}\) 上的左向量空间 (A, II, § 1.1)，而且 K 的拓扑与除环 \(K^{0}\) 的结构相容。由于这个原因，我们通常只考虑左拓扑向量空间；当我们不带限定地说“拓扑向量空间”时，应理解为指的是左向量空间。

若 \(K'\) 是 K 的子除环，而 E 是 K 上的拓扑向量空间，则显然 E 的拓扑仍然与把标量域限制到 \(K'\) 后 E 相对于 \(K'\) 的向量空间结构相容；我们说按此程序得到的 \(K'\) 上的拓扑向量空间是 K 上拓扑向量空间 E 的底空间。

为使拓扑向量空间 E 是 Hausdorff 的，必须且只需：对 E 中每个 \(x \neq 0\)，存在一个不含 x 的 0 的邻域 (GT, III, § 1.2)。

在拓扑除环 K 上的向量空间 E 上，考虑一个与其加法群结构相容的拓扑。由于恒等式

\[
\lambda x - \lambda_ {0} x _ {0} = (\lambda - \lambda_ {0}) x _ {0} + \lambda_ {0} (x - x _ {0}) + (\lambda - \lambda_ {0}) (x - x _ {0})
\]

公理 (EVT) 等价于由下面三条公理组成的公理组。

(EVT') 对所有 \(x_0 \in \mathrm{E}\)，映射 \(\lambda \mapsto \lambda x_0\) 在 \(\lambda = 0\) 处连续。

(EVT \(_{II}\) ) 对所有 \(\lambda_{0}\in K\)，映射 \(x\mapsto\lambda_{0}x\) 在 x=0 处连续。

\((\mathrm{EVT}_{\mathrm{III}}^{\prime})\) 映射 \((\lambda, x) \mapsto \lambda x\) 在 \((0, 0)\) 处连续。

特别地：

命题 1. --- 对所有 \(\alpha \in \mathbf{K}\) 及每一点 \(b \in \mathbf{E}\)，映射 \(x \mapsto \alpha x + b\)（\(\mathbf{E}\) 到自身的映射）连续。此外，若 \(\alpha \neq 0\)，则该映射是 \(\mathbf{E}\) 到自身的一个同胚。

命题的第二部分来自如下事实：若 \(\alpha \neq 0\)，则 \(x \mapsto \alpha^{-1}x - \alpha^{-1}b\) 是 \(x \mapsto \alpha x + b\) 的逆映射。

推论. --- 若 A 是 E 中的开集（相应地，闭集），则 \(\alpha A\) 对 K 中每个 \(\alpha \neq 0\) 都是 E 中的开集（相应地，闭集）。

设 E 和 F 是同一拓扑除环 K 上的两个拓扑向量空间。E 到 F 上的双射 f 是拓扑向量空间 E 到拓扑向量空间 F 的同构，当且仅当 f 是线性的且双连续。特别地，若 \(\gamma \neq 0\) 属于 K 的中心，则位似 \(x \mapsto \gamma x\) 是 E 的拓扑向量空间结构的一个自同构。

